\section{Terminal moments and the sharp constant}
\label{sec:moments}

We finish by converting the query estimates into moments and then
removing the conditioning on the prefix. Both steps use the uniformity
of all preceding bounds.

Fix a good prefix, and choose $k\in\{1,2\}$ oriented root edges
independently and uniformly with replacement from the $2m$ oriented
edges of $G$. The orientations specify the formal roots and do not
affect the query outputs. In the finite-priority model, each query at
threshold one tests survival of its root edge in the same terminal
graph. Conditional on that graph, its surviving edge fraction is
$F_n/m$, and the sampled roots are independent. Consequently
\begin{equation}\label{eq:finite-moments}
 \Prb(\text{all $k$ finite root queries return true}\mid\text{prefix})
       =\E\left[\left(\frac{F_n}{m}\right)^k
                        \mathrel{\Big|}\text{prefix}\right].
\end{equation}
This identity includes the possibility of repeated sampled edges.
Removing such samples would replace the ordinary second moment by a
different expression.

In the independent unfolding, the trees of the $k$ roots are independent,
even when their edge types coincide. Averaging over the roots gives
\begin{equation}\label{eq:unfolded-moments}
 \Prb(\text{all $k$ independent root queries return true}\mid\text{prefix})
       =\left(\frac1m\sum_{e\in\cE}q_e(1)\right)^k
       =(1+o(1))q(1)^k,
\end{equation}
by Proposition~\ref{prop:unfolding}. Lemma~\ref{lem:coupling} and
Proposition~\ref{prop:collision} bound the difference between
\eqref{eq:finite-moments} and \eqref{eq:unfolded-moments} by
\[
 \Delta_n=O(D^{-35}+n^{-1}).
\]
Since $q(1)^2=(1+2D)^{-1}$,
\begin{equation}\label{eq:relative-error}
 \frac{\Delta_n}{q(1)^2}
    =O(D^{-34}+D/n)=o(1).
\end{equation}
This also controls $\Delta_n/q(1)$, as $q(1)\le1$.

Set $Y_n=F_n/(mq(1))$, using the deterministic $m$ and $D$ from
\eqref{eq:prefix-time}. The two moment identities give, uniformly over
all good prefixes,
\[
 \E[Y_n\mid\text{prefix}]=1+o(1),\qquad
 \E[Y_n^2\mid\text{prefix}]=1+o(1).
\]
Subtracting twice the first identity from the second and adding one
yields
\begin{equation}\label{eq:conditional-l2}
 \E[(Y_n-1)^2\mid\text{prefix}]=o(1)
 \quad\text{uniformly on }\mathcal G_n.
\end{equation}

The normalization has the exact form
\begin{equation}\label{eq:constant}
 a_n:=\frac{mq(1)}{n^{3/2}}
     =\frac{\sqrt D}{2\sqrt{1+2D}}
     =\frac1{2\sqrt{2+1/D}}
     \longrightarrow a:=\frac1{2\sqrt2}.
\end{equation}
In particular $a_n$ is bounded below by a positive constant for large
$n$. Regardless of the prefix, $0\le F_n\le\binom n2\le n^2/2$.
It follows that $(Y_n-1)^2\le Cn$ deterministically. The bad-prefix
contribution therefore satisfies
\[
 \E[(Y_n-1)^2\mathbf1_{\mathcal G_n^c}]
      \le Cn\Prb(\mathcal G_n^c)=o(1).
\]
On $\mathcal G_n$, condition on the complete prefix and use
\eqref{eq:conditional-l2}. We conclude that
\[
 \E[(Y_n-1)^2]\longrightarrow0.
\]
Finally,
\[
 \E\left[\left(\frac{F_n}{n^{3/2}}-a\right)^2\right]
 \le2a_n^2\E[(Y_n-1)^2]+2(a_n-a)^2\longrightarrow0.
\]
Markov's inequality gives the convergence in probability in
Theorem~\ref{thm:main}; Cauchy--Schwarz gives the stated expectation
limit. This completes the proof.
